% ─── Lecture 24: Online Learning and Experts ─────────────────────────────────
% To include in main.tex: \section{Lecture 24: Online Learning and Experts}
%                         % ─── Lecture 24: Online Learning and Experts ─────────────────────────────────
% To include in main.tex: \section{Lecture 24: Online Learning and Experts}
%                         % ─── Lecture 24: Online Learning and Experts ─────────────────────────────────
% To include in main.tex: \section{Lecture 24: Online Learning and Experts}
%                         % ─── Lecture 24: Online Learning and Experts ─────────────────────────────────
% To include in main.tex: \section{Lecture 24: Online Learning and Experts}
%                         \input{Lecture/Lecture24}

\subsection{Introduction: Online vs Offline Algorithms}

\begin{definition}{Offline Algorithm}{offline}
An algorithm where the \textbf{entire input is known in advance}. The algorithm can plan with full information.\\[0.3em]
\textit{Examples:} sorting an array, finding a Minimum Spanning Tree (MST).
\end{definition}

\begin{definition}{Online Algorithm}{online}
An algorithm where \textbf{input arrives one piece at a time}. The algorithm makes \textbf{irrevocable decisions} without knowledge of future requests.\\[0.3em]
\textit{Key principle:} The algorithm must act before it knows the whole story.
\end{definition}

Online decision problems arise in many practical settings:

\begin{center}
\begin{tabular}{ll}
\toprule
\textbf{Setting} & \textbf{Online Decision} \\
\midrule
Cache eviction    & Which page to evict? \\
Online ads        & Which ad to show? \\
Rent or buy       & Keep renting or buy? \\
Expert advice     & Whose advice today? \\
Hiring/selection  & Accept now or wait? \\
\bottomrule
\end{tabular}
\end{center}

A reasonable decision may look poor after the future is revealed. The relevant question is not "was the algorithm psychic?" but rather: \textbf{how badly can it be punished for not being psychic?}

% ─────────────────────────────────────────────────────────────────────────────
\subsection{Measuring Performance}

\begin{definition}{Competitive Ratio}{competitive}
Compare the algorithm's cost with the best solution that knew the whole input in advance (the offline optimum):
\[
\mathrm{ALG}(I) \leq r \cdot \mathrm{OPT}(I) + c
\]
Used when the offline optimum is a meaningful and stable target.
\end{definition}

When comparing to the full offline optimum is too ambitious, we use a weaker but still meaningful benchmark.

\begin{definition}{Regret}{regret}
Compare the algorithm's total loss with the \textbf{best single expert in hindsight}. The extra loss incurred by the algorithm is called \textbf{regret}:
\[
\text{regret} = \mathrm{loss}(\mathrm{ALG}) - \min_i \mathrm{loss}(i)
\]
Used when "best possible" is too strong or too unstable as a benchmark.
\end{definition}

\begin{remark}{Performance benchmarks}{bench}
\textbf{Happiness} $\approx$ outcome $-$ expectation, but happiness is also \textit{relative} to the people next to you. Similarly, algorithm performance is judged relative to what was achievable.
\end{remark}

% ─────────────────────────────────────────────────────────────────────────────
\subsection{Prediction with Expert Advice}

\subsubsection{The Expert Setting}

Consider the following scenario: every morning, $N$ experts each predict YES or NO (e.g., "will it rain?"). We must make our own prediction before observing the outcome.

\textbf{The game.} For $t = 1, \ldots, T$:
\begin{enumerate}
  \item Each expert $i \in \{1, \ldots, N\}$ predicts YES or NO.
  \item The algorithm predicts YES or NO.
  \item The true outcome is revealed.
  \item Everyone who predicted incorrectly receives one \textit{mistake}.
\end{enumerate}

Let $M_{\mathrm{ALG}}$ denote the number of mistakes of the algorithm, and $M_i$ the number of mistakes of expert $i$.

\textbf{Goal:} Keep $M_{\mathrm{ALG}}$ close to $\min_i M_i$.

\subsubsection{Why Compare to the Best Expert?}

\begin{itemize}
  \item We cannot promise to make few mistakes on every sequence (e.g., if all experts say YES and the answer is NO, any algorithm using their advice is in trouble).
  \item But if one expert was consistently good, we would like to discover this quickly.
  \item \textbf{Regret} asks: is $M_{\mathrm{ALG}} - \min_i M_i$ small?
\end{itemize}

% ─────────────────────────────────────────────────────────────────────────────
\subsection{Warmup: One Perfect Expert}

\subsubsection{The Elimination Algorithm}

Suppose at least one expert makes \emph{zero} mistakes. We can exploit this with a simple elimination strategy.

\begin{mdframed}[backgroundcolor=lightblue, linecolor=keyblue, linewidth=1pt]
\textbf{Elimination Algorithm}
\begin{enumerate}
  \item Keep the experts who have made \textbf{no mistakes} so far.
  \item Predict by \textbf{majority vote} among the remaining experts.
  \item After the outcome is revealed, \textbf{delete} every expert who was wrong.
\end{enumerate}
\end{mdframed}

\textbf{Correctness:} The perfect expert is never deleted, so the set of remaining experts is never empty.

\subsubsection{Analysis}

\begin{theorem}{Perfect Expert Guarantee}{perfect}
If at least one expert makes no mistakes, the Elimination algorithm makes \textbf{at most $\lg N$ mistakes}.
\end{theorem}

\textbf{Proof sketch.} Suppose the algorithm makes a mistake. Since it followed the majority of remaining experts, at least half of the remaining experts were wrong — and are now deleted. The pool size decreases:
\[
N \;\to\; \frac{N}{2} \;\to\; \frac{N}{4} \;\to\; \cdots
\]
After more than $\lg N$ mistakes, fewer than one expert would remain — impossible since the perfect expert survives. \hfill$\square$

\textbf{Example:} With $N = 1024$ experts, even a very adversarial sequence causes at most 10 mistakes before the algorithm has essentially identified the perfect expert.

\subsubsection{When Nobody is Perfect}

A first idea: run Elimination until all experts are deleted, then restart with all experts alive. If the best expert makes $M$ mistakes, there are at most $M+1$ restarts, giving:
\[
M_{\mathrm{ALG}} \leq (M+1) \lg N
\]
This bound is often \textbf{too weak}: the overhead grows \emph{multiplicatively} with $M$. We want an \textbf{additive} penalty instead.

% ─────────────────────────────────────────────────────────────────────────────
\subsection{Weighted Majority}

\subsubsection{Soft Elimination}

Instead of deleting an expert after a mistake, we \emph{reduce its influence}:
\begin{itemize}
  \item If correct today: \textbf{keep weight}.
  \item If wrong today: \textbf{halve weight}.
  \item Wrong often: weight \textbf{becomes negligible}.
\end{itemize}
The algorithm listens to a \textbf{weighted majority vote}.

\subsubsection{The Algorithm}

\begin{mdframed}[backgroundcolor=lightblue, linecolor=keyblue, linewidth=1pt]
\textbf{Weighted Majority Algorithm}

Initially set $w_i^{(1)} = 1$ for every expert $i$. For $t = 1, \ldots, T$:
\begin{enumerate}
  \item Each expert predicts YES or NO.
  \item Predict the answer with \textbf{larger total weight}.
  \item After the outcome is revealed, update:
    \[
      w_i^{(t+1)} = \begin{cases} w_i^{(t)} / 2 & \text{if expert } i \text{ was wrong,} \\ w_i^{(t)} & \text{otherwise.} \end{cases}
    \]
\end{enumerate}
\end{mdframed}

Weights act as \textbf{trust scores}: they only decrease when the expert makes a mistake, so a good expert retains much more weight than a bad one.

\subsubsection{Guarantee}

\begin{theorem}{Weighted Majority Bound}{wm}
Let $M_{\mathrm{WM}}$ be the number of mistakes of Weighted Majority. For every expert $i$:
\[
M_{\mathrm{WM}} \leq 2.41 \cdot (M_i + \lg N)
\]
Therefore:
\[
M_{\mathrm{WM}} \leq 2.41 \cdot \left(\min_i M_i + \lg N\right)
\]
\end{theorem}

This is much better than the restarting bound when $M_i$ is large: the overhead is now \textbf{additive} in $\lg N$.

\subsubsection{Analysis via Potential Function}

Define the \textbf{potential} (total weight):
\[
\Phi^{(t)} = \sum_{i=1}^{N} w_i^{(t)}
\]

\textbf{Upper bound} (from algorithm's mistakes): When the algorithm makes a mistake, it followed the weighted majority, so at least half the total weight was on wrong experts. Those weights are halved, giving:
\[
\Phi^{(t+1)} \leq \frac{1}{2}\Phi^{(t)} + \frac{1}{2} \cdot \frac{1}{2}\Phi^{(t)} = \frac{3}{4}\Phi^{(t)}
\]
After $M_{\mathrm{WM}}$ mistakes:
\[
\Phi^{(T+1)} \leq N \left(\frac{3}{4}\right)^{M_{\mathrm{WM}}}
\]

\textbf{Lower bound} (from a good expert): If expert $i$ makes $M_i$ mistakes, its final weight is $(1/2)^{M_i}$. Since the total weight is at least this:
\[
\left(\frac{1}{2}\right)^{M_i} \leq \Phi^{(T+1)}
\]

\textbf{Combining:}
\[
\left(\frac{1}{2}\right)^{M_i} \leq N \left(\frac{3}{4}\right)^{M_{\mathrm{WM}}}
\]
Taking base-2 logarithms and rearranging:
\[
M_{\mathrm{WM}} \leq \frac{1}{\lg(4/3)}(M_i + \lg N) \approx 2.41\,(M_i + \lg N)
\]

\textbf{Intuition:} Many algorithm mistakes $\Rightarrow$ total weight becomes tiny. A good expert prevents this from happening.

\subsubsection{Tuning the Penalty}

Halving (multiplying by $1/2$ on a mistake) is a special case. More generally, we can multiply by $(1-\varepsilon)$ for any $\varepsilon \in (0,1)$.

\begin{theorem}{Tunable Weighted Majority}{tunable}
With a suitable weighted-majority update using parameter $\varepsilon$:
\[
M_{\mathrm{WM}} \leq 2(1+\varepsilon)\,M_i + O\!\left(\frac{\lg N}{\varepsilon}\right) \quad \text{for every expert } i
\]
\end{theorem}

\textbf{Tradeoff:} Smaller $\varepsilon$ means less overreaction to mistakes (more patient), but slower learning.

% ─────────────────────────────────────────────────────────────────────────────
\subsection{Randomized Experts: The Hedge Algorithm}

\subsubsection{Limitation of Weighted Majority}

Weighted Majority is \textbf{deterministic}. An adversary who can observe our prediction can always choose the opposite outcome, forcing a mistake every round. In fact, there exist sequences where deterministic Weighted Majority makes \emph{about twice} as many mistakes as the best expert. To improve the constant, we stop committing to one answer deterministically.

\subsubsection{Randomized Prediction}

Instead of choosing one expert's advice deterministically, we choose a \textbf{distribution} $\mathbf{p}^{(t)} = (p_1^{(t)}, \ldots, p_N^{(t)})$ over experts. If expert $i$ incurs loss $m_i^{(t)} \in [-1, 1]$ on day $t$, our \textbf{expected loss} is:
\[
\sum_{i=1}^N p_i^{(t)} m_i^{(t)} = \mathbf{p}^{(t)} \cdot \mathbf{m}^{(t)}
\]
This framework covers binary mistakes, but also richer cost structures: money lost, error magnitude, reward minus cost, etc.

\subsubsection{The Hedge Algorithm}

\begin{mdframed}[backgroundcolor=lightblue, linecolor=keyblue, linewidth=1pt]
\textbf{Hedge Algorithm} (parameter $\varepsilon > 0$)

Maintain weights $w_i^{(t)}$. At time $t$:
\begin{enumerate}
  \item Play expert $i$ with probability
    \[
      p_i^{(t)} = \frac{w_i^{(t)}}{\Phi^{(t)}}, \quad \text{where } \Phi^{(t)} = \sum_j w_j^{(t)}
    \]
  \item Observe the loss vector $\mathbf{m}^{(t)}$.
  \item Update weights:
    \[
      w_i^{(t+1)} = w_i^{(t)} \cdot e^{-\varepsilon m_i^{(t)}}
    \]
\end{enumerate}
\end{mdframed}

Low loss $\Rightarrow$ weight increases; high loss $\Rightarrow$ weight decreases. The key distinction from Weighted Majority is the use of the \textbf{exponential update}, which handles continuous losses naturally.

\subsubsection{Hedge Guarantee}

\begin{theorem}{Hedge Regret Bound}{hedge}
If $m_i^{(t)} \in [-1,1]$ and $\varepsilon \leq 1$, then for every expert $i$:
\[
\sum_{t=1}^T \mathbf{p}^{(t)} \cdot \mathbf{m}^{(t)} \;\leq\; \sum_{t=1}^T m_i^{(t)} + \frac{\ln N}{\varepsilon} + \varepsilon T
\]
\end{theorem}

Choosing $\varepsilon \approx \sqrt{\ln N / T}$ balances the two additive terms and gives:
\[
\text{regret} = O\!\left(\sqrt{T \ln N}\right)
\]
The \textbf{average regret} vanishes as $T \to \infty$:
\[
\frac{\text{regret}}{T} = O\!\left(\sqrt{\frac{\ln N}{T}}\right) \;\to\; 0
\]
This is called \textbf{no-regret} or \textbf{Hannan consistency}.

% ─────────────────────────────────────────────────────────────────────────────
\subsection{Why Multiplicative Weights is Everywhere}

The Multiplicative Weights / Hedge framework is a small idea with an enormous footprint across computer science:

\begin{itemize}
  \item \textbf{Boosting} (ML): combine many weak classifiers into a strong one (AdaBoost is a special case).
  \item \textbf{Game theory}: repeated play with Hedge converges to correlated equilibria.
  \item \textbf{Optimization}: fast approximate algorithms for linear programs and semidefinite programs (SDPs).
  \item \textbf{Online learning}: portfolio selection, prediction markets, and adaptive decision making.
  \item \textbf{Algorithm design}: a standard potential-function proof template.
\end{itemize}

\begin{remark}{The recurring pattern}{mw}
\textbf{Multiplicatively reward what worked, but do not completely forget alternatives.}\\[0.3em]
This is the fundamental principle behind multiplicative weights: unlike additive updates that can send weights negative or zero, exponential updates keep all options alive while shifting mass toward better-performing experts.
\end{remark}

% ─────────────────────────────────────────────────────────────────────────────
\subsection{Summary}

\begin{center}
\rowcolors{2}{gray!10}{white}
\begin{tabular}{p{3cm} p{4.5cm} p{5cm}}
\toprule
\textbf{Algorithm} & \textbf{Guarantee} & \textbf{Setting} \\
\midrule
Elimination & $M_{\mathrm{ALG}} \leq \lg N$ & One perfect expert exists \\
Restart Elimination & $M_{\mathrm{ALG}} \leq (M^*+1)\lg N$ & General, but multiplicative overhead \\
Weighted Majority & $M_{\mathrm{WM}} \leq 2.41(M^* + \lg N)$ & General, deterministic \\
Hedge & $\text{regret} = O(\sqrt{T\ln N})$ & General, randomized, continuous losses \\
\bottomrule
\end{tabular}
\end{center}

where $M^* = \min_i M_i$ is the number of mistakes of the best expert.

\textbf{Key takeaways:}
\begin{itemize}
  \item Online algorithms must make irrevocable decisions without seeing the future.
  \item When comparing to full hindsight is too ambitious, \textbf{regret} is the right benchmark.
  \item The \textbf{Elimination algorithm} achieves $\lg N$ mistakes with a perfect expert.
  \item \textbf{Weighted Majority} handles imperfect experts with an additive $O(\lg N)$ overhead, analyzed via a potential function argument.
  \item \textbf{Hedge} uses randomization and exponential updates to achieve sublinear $O(\sqrt{T \ln N})$ regret.
\end{itemize}

\subsection{Introduction: Online vs Offline Algorithms}

\begin{definition}{Offline Algorithm}{offline}
An algorithm where the \textbf{entire input is known in advance}. The algorithm can plan with full information.\\[0.3em]
\textit{Examples:} sorting an array, finding a Minimum Spanning Tree (MST).
\end{definition}

\begin{definition}{Online Algorithm}{online}
An algorithm where \textbf{input arrives one piece at a time}. The algorithm makes \textbf{irrevocable decisions} without knowledge of future requests.\\[0.3em]
\textit{Key principle:} The algorithm must act before it knows the whole story.
\end{definition}

Online decision problems arise in many practical settings:

\begin{center}
\begin{tabular}{ll}
\toprule
\textbf{Setting} & \textbf{Online Decision} \\
\midrule
Cache eviction    & Which page to evict? \\
Online ads        & Which ad to show? \\
Rent or buy       & Keep renting or buy? \\
Expert advice     & Whose advice today? \\
Hiring/selection  & Accept now or wait? \\
\bottomrule
\end{tabular}
\end{center}

A reasonable decision may look poor after the future is revealed. The relevant question is not "was the algorithm psychic?" but rather: \textbf{how badly can it be punished for not being psychic?}

% ─────────────────────────────────────────────────────────────────────────────
\subsection{Measuring Performance}

\begin{definition}{Competitive Ratio}{competitive}
Compare the algorithm's cost with the best solution that knew the whole input in advance (the offline optimum):
\[
\mathrm{ALG}(I) \leq r \cdot \mathrm{OPT}(I) + c
\]
Used when the offline optimum is a meaningful and stable target.
\end{definition}

When comparing to the full offline optimum is too ambitious, we use a weaker but still meaningful benchmark.

\begin{definition}{Regret}{regret}
Compare the algorithm's total loss with the \textbf{best single expert in hindsight}. The extra loss incurred by the algorithm is called \textbf{regret}:
\[
\text{regret} = \mathrm{loss}(\mathrm{ALG}) - \min_i \mathrm{loss}(i)
\]
Used when "best possible" is too strong or too unstable as a benchmark.
\end{definition}

\begin{remark}{Performance benchmarks}{bench}
\textbf{Happiness} $\approx$ outcome $-$ expectation, but happiness is also \textit{relative} to the people next to you. Similarly, algorithm performance is judged relative to what was achievable.
\end{remark}

% ─────────────────────────────────────────────────────────────────────────────
\subsection{Prediction with Expert Advice}

\subsubsection{The Expert Setting}

Consider the following scenario: every morning, $N$ experts each predict YES or NO (e.g., "will it rain?"). We must make our own prediction before observing the outcome.

\textbf{The game.} For $t = 1, \ldots, T$:
\begin{enumerate}
  \item Each expert $i \in \{1, \ldots, N\}$ predicts YES or NO.
  \item The algorithm predicts YES or NO.
  \item The true outcome is revealed.
  \item Everyone who predicted incorrectly receives one \textit{mistake}.
\end{enumerate}

Let $M_{\mathrm{ALG}}$ denote the number of mistakes of the algorithm, and $M_i$ the number of mistakes of expert $i$.

\textbf{Goal:} Keep $M_{\mathrm{ALG}}$ close to $\min_i M_i$.

\subsubsection{Why Compare to the Best Expert?}

\begin{itemize}
  \item We cannot promise to make few mistakes on every sequence (e.g., if all experts say YES and the answer is NO, any algorithm using their advice is in trouble).
  \item But if one expert was consistently good, we would like to discover this quickly.
  \item \textbf{Regret} asks: is $M_{\mathrm{ALG}} - \min_i M_i$ small?
\end{itemize}

% ─────────────────────────────────────────────────────────────────────────────
\subsection{Warmup: One Perfect Expert}

\subsubsection{The Elimination Algorithm}

Suppose at least one expert makes \emph{zero} mistakes. We can exploit this with a simple elimination strategy.

\begin{mdframed}[backgroundcolor=lightblue, linecolor=keyblue, linewidth=1pt]
\textbf{Elimination Algorithm}
\begin{enumerate}
  \item Keep the experts who have made \textbf{no mistakes} so far.
  \item Predict by \textbf{majority vote} among the remaining experts.
  \item After the outcome is revealed, \textbf{delete} every expert who was wrong.
\end{enumerate}
\end{mdframed}

\textbf{Correctness:} The perfect expert is never deleted, so the set of remaining experts is never empty.

\subsubsection{Analysis}

\begin{theorem}{Perfect Expert Guarantee}{perfect}
If at least one expert makes no mistakes, the Elimination algorithm makes \textbf{at most $\lg N$ mistakes}.
\end{theorem}

\textbf{Proof sketch.} Suppose the algorithm makes a mistake. Since it followed the majority of remaining experts, at least half of the remaining experts were wrong — and are now deleted. The pool size decreases:
\[
N \;\to\; \frac{N}{2} \;\to\; \frac{N}{4} \;\to\; \cdots
\]
After more than $\lg N$ mistakes, fewer than one expert would remain — impossible since the perfect expert survives. \hfill$\square$

\textbf{Example:} With $N = 1024$ experts, even a very adversarial sequence causes at most 10 mistakes before the algorithm has essentially identified the perfect expert.

\subsubsection{When Nobody is Perfect}

A first idea: run Elimination until all experts are deleted, then restart with all experts alive. If the best expert makes $M$ mistakes, there are at most $M+1$ restarts, giving:
\[
M_{\mathrm{ALG}} \leq (M+1) \lg N
\]
This bound is often \textbf{too weak}: the overhead grows \emph{multiplicatively} with $M$. We want an \textbf{additive} penalty instead.

% ─────────────────────────────────────────────────────────────────────────────
\subsection{Weighted Majority}

\subsubsection{Soft Elimination}

Instead of deleting an expert after a mistake, we \emph{reduce its influence}:
\begin{itemize}
  \item If correct today: \textbf{keep weight}.
  \item If wrong today: \textbf{halve weight}.
  \item Wrong often: weight \textbf{becomes negligible}.
\end{itemize}
The algorithm listens to a \textbf{weighted majority vote}.

\subsubsection{The Algorithm}

\begin{mdframed}[backgroundcolor=lightblue, linecolor=keyblue, linewidth=1pt]
\textbf{Weighted Majority Algorithm}

Initially set $w_i^{(1)} = 1$ for every expert $i$. For $t = 1, \ldots, T$:
\begin{enumerate}
  \item Each expert predicts YES or NO.
  \item Predict the answer with \textbf{larger total weight}.
  \item After the outcome is revealed, update:
    \[
      w_i^{(t+1)} = \begin{cases} w_i^{(t)} / 2 & \text{if expert } i \text{ was wrong,} \\ w_i^{(t)} & \text{otherwise.} \end{cases}
    \]
\end{enumerate}
\end{mdframed}

Weights act as \textbf{trust scores}: they only decrease when the expert makes a mistake, so a good expert retains much more weight than a bad one.

\subsubsection{Guarantee}

\begin{theorem}{Weighted Majority Bound}{wm}
Let $M_{\mathrm{WM}}$ be the number of mistakes of Weighted Majority. For every expert $i$:
\[
M_{\mathrm{WM}} \leq 2.41 \cdot (M_i + \lg N)
\]
Therefore:
\[
M_{\mathrm{WM}} \leq 2.41 \cdot \left(\min_i M_i + \lg N\right)
\]
\end{theorem}

This is much better than the restarting bound when $M_i$ is large: the overhead is now \textbf{additive} in $\lg N$.

\subsubsection{Analysis via Potential Function}

Define the \textbf{potential} (total weight):
\[
\Phi^{(t)} = \sum_{i=1}^{N} w_i^{(t)}
\]

\textbf{Upper bound} (from algorithm's mistakes): When the algorithm makes a mistake, it followed the weighted majority, so at least half the total weight was on wrong experts. Those weights are halved, giving:
\[
\Phi^{(t+1)} \leq \frac{1}{2}\Phi^{(t)} + \frac{1}{2} \cdot \frac{1}{2}\Phi^{(t)} = \frac{3}{4}\Phi^{(t)}
\]
After $M_{\mathrm{WM}}$ mistakes:
\[
\Phi^{(T+1)} \leq N \left(\frac{3}{4}\right)^{M_{\mathrm{WM}}}
\]

\textbf{Lower bound} (from a good expert): If expert $i$ makes $M_i$ mistakes, its final weight is $(1/2)^{M_i}$. Since the total weight is at least this:
\[
\left(\frac{1}{2}\right)^{M_i} \leq \Phi^{(T+1)}
\]

\textbf{Combining:}
\[
\left(\frac{1}{2}\right)^{M_i} \leq N \left(\frac{3}{4}\right)^{M_{\mathrm{WM}}}
\]
Taking base-2 logarithms and rearranging:
\[
M_{\mathrm{WM}} \leq \frac{1}{\lg(4/3)}(M_i + \lg N) \approx 2.41\,(M_i + \lg N)
\]

\textbf{Intuition:} Many algorithm mistakes $\Rightarrow$ total weight becomes tiny. A good expert prevents this from happening.

\subsubsection{Tuning the Penalty}

Halving (multiplying by $1/2$ on a mistake) is a special case. More generally, we can multiply by $(1-\varepsilon)$ for any $\varepsilon \in (0,1)$.

\begin{theorem}{Tunable Weighted Majority}{tunable}
With a suitable weighted-majority update using parameter $\varepsilon$:
\[
M_{\mathrm{WM}} \leq 2(1+\varepsilon)\,M_i + O\!\left(\frac{\lg N}{\varepsilon}\right) \quad \text{for every expert } i
\]
\end{theorem}

\textbf{Tradeoff:} Smaller $\varepsilon$ means less overreaction to mistakes (more patient), but slower learning.

% ─────────────────────────────────────────────────────────────────────────────
\subsection{Randomized Experts: The Hedge Algorithm}

\subsubsection{Limitation of Weighted Majority}

Weighted Majority is \textbf{deterministic}. An adversary who can observe our prediction can always choose the opposite outcome, forcing a mistake every round. In fact, there exist sequences where deterministic Weighted Majority makes \emph{about twice} as many mistakes as the best expert. To improve the constant, we stop committing to one answer deterministically.

\subsubsection{Randomized Prediction}

Instead of choosing one expert's advice deterministically, we choose a \textbf{distribution} $\mathbf{p}^{(t)} = (p_1^{(t)}, \ldots, p_N^{(t)})$ over experts. If expert $i$ incurs loss $m_i^{(t)} \in [-1, 1]$ on day $t$, our \textbf{expected loss} is:
\[
\sum_{i=1}^N p_i^{(t)} m_i^{(t)} = \mathbf{p}^{(t)} \cdot \mathbf{m}^{(t)}
\]
This framework covers binary mistakes, but also richer cost structures: money lost, error magnitude, reward minus cost, etc.

\subsubsection{The Hedge Algorithm}

\begin{mdframed}[backgroundcolor=lightblue, linecolor=keyblue, linewidth=1pt]
\textbf{Hedge Algorithm} (parameter $\varepsilon > 0$)

Maintain weights $w_i^{(t)}$. At time $t$:
\begin{enumerate}
  \item Play expert $i$ with probability
    \[
      p_i^{(t)} = \frac{w_i^{(t)}}{\Phi^{(t)}}, \quad \text{where } \Phi^{(t)} = \sum_j w_j^{(t)}
    \]
  \item Observe the loss vector $\mathbf{m}^{(t)}$.
  \item Update weights:
    \[
      w_i^{(t+1)} = w_i^{(t)} \cdot e^{-\varepsilon m_i^{(t)}}
    \]
\end{enumerate}
\end{mdframed}

Low loss $\Rightarrow$ weight increases; high loss $\Rightarrow$ weight decreases. The key distinction from Weighted Majority is the use of the \textbf{exponential update}, which handles continuous losses naturally.

\subsubsection{Hedge Guarantee}

\begin{theorem}{Hedge Regret Bound}{hedge}
If $m_i^{(t)} \in [-1,1]$ and $\varepsilon \leq 1$, then for every expert $i$:
\[
\sum_{t=1}^T \mathbf{p}^{(t)} \cdot \mathbf{m}^{(t)} \;\leq\; \sum_{t=1}^T m_i^{(t)} + \frac{\ln N}{\varepsilon} + \varepsilon T
\]
\end{theorem}

Choosing $\varepsilon \approx \sqrt{\ln N / T}$ balances the two additive terms and gives:
\[
\text{regret} = O\!\left(\sqrt{T \ln N}\right)
\]
The \textbf{average regret} vanishes as $T \to \infty$:
\[
\frac{\text{regret}}{T} = O\!\left(\sqrt{\frac{\ln N}{T}}\right) \;\to\; 0
\]
This is called \textbf{no-regret} or \textbf{Hannan consistency}.

% ─────────────────────────────────────────────────────────────────────────────
\subsection{Why Multiplicative Weights is Everywhere}

The Multiplicative Weights / Hedge framework is a small idea with an enormous footprint across computer science:

\begin{itemize}
  \item \textbf{Boosting} (ML): combine many weak classifiers into a strong one (AdaBoost is a special case).
  \item \textbf{Game theory}: repeated play with Hedge converges to correlated equilibria.
  \item \textbf{Optimization}: fast approximate algorithms for linear programs and semidefinite programs (SDPs).
  \item \textbf{Online learning}: portfolio selection, prediction markets, and adaptive decision making.
  \item \textbf{Algorithm design}: a standard potential-function proof template.
\end{itemize}

\begin{remark}{The recurring pattern}{mw}
\textbf{Multiplicatively reward what worked, but do not completely forget alternatives.}\\[0.3em]
This is the fundamental principle behind multiplicative weights: unlike additive updates that can send weights negative or zero, exponential updates keep all options alive while shifting mass toward better-performing experts.
\end{remark}

% ─────────────────────────────────────────────────────────────────────────────
\subsection{Summary}

\begin{center}
\rowcolors{2}{gray!10}{white}
\begin{tabular}{p{3cm} p{4.5cm} p{5cm}}
\toprule
\textbf{Algorithm} & \textbf{Guarantee} & \textbf{Setting} \\
\midrule
Elimination & $M_{\mathrm{ALG}} \leq \lg N$ & One perfect expert exists \\
Restart Elimination & $M_{\mathrm{ALG}} \leq (M^*+1)\lg N$ & General, but multiplicative overhead \\
Weighted Majority & $M_{\mathrm{WM}} \leq 2.41(M^* + \lg N)$ & General, deterministic \\
Hedge & $\text{regret} = O(\sqrt{T\ln N})$ & General, randomized, continuous losses \\
\bottomrule
\end{tabular}
\end{center}

where $M^* = \min_i M_i$ is the number of mistakes of the best expert.

\textbf{Key takeaways:}
\begin{itemize}
  \item Online algorithms must make irrevocable decisions without seeing the future.
  \item When comparing to full hindsight is too ambitious, \textbf{regret} is the right benchmark.
  \item The \textbf{Elimination algorithm} achieves $\lg N$ mistakes with a perfect expert.
  \item \textbf{Weighted Majority} handles imperfect experts with an additive $O(\lg N)$ overhead, analyzed via a potential function argument.
  \item \textbf{Hedge} uses randomization and exponential updates to achieve sublinear $O(\sqrt{T \ln N})$ regret.
\end{itemize}

\subsection{Introduction: Online vs Offline Algorithms}

\begin{definition}{Offline Algorithm}{offline}
An algorithm where the \textbf{entire input is known in advance}. The algorithm can plan with full information.\\[0.3em]
\textit{Examples:} sorting an array, finding a Minimum Spanning Tree (MST).
\end{definition}

\begin{definition}{Online Algorithm}{online}
An algorithm where \textbf{input arrives one piece at a time}. The algorithm makes \textbf{irrevocable decisions} without knowledge of future requests.\\[0.3em]
\textit{Key principle:} The algorithm must act before it knows the whole story.
\end{definition}

Online decision problems arise in many practical settings:

\begin{center}
\begin{tabular}{ll}
\toprule
\textbf{Setting} & \textbf{Online Decision} \\
\midrule
Cache eviction    & Which page to evict? \\
Online ads        & Which ad to show? \\
Rent or buy       & Keep renting or buy? \\
Expert advice     & Whose advice today? \\
Hiring/selection  & Accept now or wait? \\
\bottomrule
\end{tabular}
\end{center}

A reasonable decision may look poor after the future is revealed. The relevant question is not "was the algorithm psychic?" but rather: \textbf{how badly can it be punished for not being psychic?}

% ─────────────────────────────────────────────────────────────────────────────
\subsection{Measuring Performance}

\begin{definition}{Competitive Ratio}{competitive}
Compare the algorithm's cost with the best solution that knew the whole input in advance (the offline optimum):
\[
\mathrm{ALG}(I) \leq r \cdot \mathrm{OPT}(I) + c
\]
Used when the offline optimum is a meaningful and stable target.
\end{definition}

When comparing to the full offline optimum is too ambitious, we use a weaker but still meaningful benchmark.

\begin{definition}{Regret}{regret}
Compare the algorithm's total loss with the \textbf{best single expert in hindsight}. The extra loss incurred by the algorithm is called \textbf{regret}:
\[
\text{regret} = \mathrm{loss}(\mathrm{ALG}) - \min_i \mathrm{loss}(i)
\]
Used when "best possible" is too strong or too unstable as a benchmark.
\end{definition}

\begin{remark}{Performance benchmarks}{bench}
\textbf{Happiness} $\approx$ outcome $-$ expectation, but happiness is also \textit{relative} to the people next to you. Similarly, algorithm performance is judged relative to what was achievable.
\end{remark}

% ─────────────────────────────────────────────────────────────────────────────
\subsection{Prediction with Expert Advice}

\subsubsection{The Expert Setting}

Consider the following scenario: every morning, $N$ experts each predict YES or NO (e.g., "will it rain?"). We must make our own prediction before observing the outcome.

\textbf{The game.} For $t = 1, \ldots, T$:
\begin{enumerate}
  \item Each expert $i \in \{1, \ldots, N\}$ predicts YES or NO.
  \item The algorithm predicts YES or NO.
  \item The true outcome is revealed.
  \item Everyone who predicted incorrectly receives one \textit{mistake}.
\end{enumerate}

Let $M_{\mathrm{ALG}}$ denote the number of mistakes of the algorithm, and $M_i$ the number of mistakes of expert $i$.

\textbf{Goal:} Keep $M_{\mathrm{ALG}}$ close to $\min_i M_i$.

\subsubsection{Why Compare to the Best Expert?}

\begin{itemize}
  \item We cannot promise to make few mistakes on every sequence (e.g., if all experts say YES and the answer is NO, any algorithm using their advice is in trouble).
  \item But if one expert was consistently good, we would like to discover this quickly.
  \item \textbf{Regret} asks: is $M_{\mathrm{ALG}} - \min_i M_i$ small?
\end{itemize}

% ─────────────────────────────────────────────────────────────────────────────
\subsection{Warmup: One Perfect Expert}

\subsubsection{The Elimination Algorithm}

Suppose at least one expert makes \emph{zero} mistakes. We can exploit this with a simple elimination strategy.

\begin{mdframed}[backgroundcolor=lightblue, linecolor=keyblue, linewidth=1pt]
\textbf{Elimination Algorithm}
\begin{enumerate}
  \item Keep the experts who have made \textbf{no mistakes} so far.
  \item Predict by \textbf{majority vote} among the remaining experts.
  \item After the outcome is revealed, \textbf{delete} every expert who was wrong.
\end{enumerate}
\end{mdframed}

\textbf{Correctness:} The perfect expert is never deleted, so the set of remaining experts is never empty.

\subsubsection{Analysis}

\begin{theorem}{Perfect Expert Guarantee}{perfect}
If at least one expert makes no mistakes, the Elimination algorithm makes \textbf{at most $\lg N$ mistakes}.
\end{theorem}

\textbf{Proof sketch.} Suppose the algorithm makes a mistake. Since it followed the majority of remaining experts, at least half of the remaining experts were wrong — and are now deleted. The pool size decreases:
\[
N \;\to\; \frac{N}{2} \;\to\; \frac{N}{4} \;\to\; \cdots
\]
After more than $\lg N$ mistakes, fewer than one expert would remain — impossible since the perfect expert survives. \hfill$\square$

\textbf{Example:} With $N = 1024$ experts, even a very adversarial sequence causes at most 10 mistakes before the algorithm has essentially identified the perfect expert.

\subsubsection{When Nobody is Perfect}

A first idea: run Elimination until all experts are deleted, then restart with all experts alive. If the best expert makes $M$ mistakes, there are at most $M+1$ restarts, giving:
\[
M_{\mathrm{ALG}} \leq (M+1) \lg N
\]
This bound is often \textbf{too weak}: the overhead grows \emph{multiplicatively} with $M$. We want an \textbf{additive} penalty instead.

% ─────────────────────────────────────────────────────────────────────────────
\subsection{Weighted Majority}

\subsubsection{Soft Elimination}

Instead of deleting an expert after a mistake, we \emph{reduce its influence}:
\begin{itemize}
  \item If correct today: \textbf{keep weight}.
  \item If wrong today: \textbf{halve weight}.
  \item Wrong often: weight \textbf{becomes negligible}.
\end{itemize}
The algorithm listens to a \textbf{weighted majority vote}.

\subsubsection{The Algorithm}

\begin{mdframed}[backgroundcolor=lightblue, linecolor=keyblue, linewidth=1pt]
\textbf{Weighted Majority Algorithm}

Initially set $w_i^{(1)} = 1$ for every expert $i$. For $t = 1, \ldots, T$:
\begin{enumerate}
  \item Each expert predicts YES or NO.
  \item Predict the answer with \textbf{larger total weight}.
  \item After the outcome is revealed, update:
    \[
      w_i^{(t+1)} = \begin{cases} w_i^{(t)} / 2 & \text{if expert } i \text{ was wrong,} \\ w_i^{(t)} & \text{otherwise.} \end{cases}
    \]
\end{enumerate}
\end{mdframed}

Weights act as \textbf{trust scores}: they only decrease when the expert makes a mistake, so a good expert retains much more weight than a bad one.

\subsubsection{Guarantee}

\begin{theorem}{Weighted Majority Bound}{wm}
Let $M_{\mathrm{WM}}$ be the number of mistakes of Weighted Majority. For every expert $i$:
\[
M_{\mathrm{WM}} \leq 2.41 \cdot (M_i + \lg N)
\]
Therefore:
\[
M_{\mathrm{WM}} \leq 2.41 \cdot \left(\min_i M_i + \lg N\right)
\]
\end{theorem}

This is much better than the restarting bound when $M_i$ is large: the overhead is now \textbf{additive} in $\lg N$.

\subsubsection{Analysis via Potential Function}

Define the \textbf{potential} (total weight):
\[
\Phi^{(t)} = \sum_{i=1}^{N} w_i^{(t)}
\]

\textbf{Upper bound} (from algorithm's mistakes): When the algorithm makes a mistake, it followed the weighted majority, so at least half the total weight was on wrong experts. Those weights are halved, giving:
\[
\Phi^{(t+1)} \leq \frac{1}{2}\Phi^{(t)} + \frac{1}{2} \cdot \frac{1}{2}\Phi^{(t)} = \frac{3}{4}\Phi^{(t)}
\]
After $M_{\mathrm{WM}}$ mistakes:
\[
\Phi^{(T+1)} \leq N \left(\frac{3}{4}\right)^{M_{\mathrm{WM}}}
\]

\textbf{Lower bound} (from a good expert): If expert $i$ makes $M_i$ mistakes, its final weight is $(1/2)^{M_i}$. Since the total weight is at least this:
\[
\left(\frac{1}{2}\right)^{M_i} \leq \Phi^{(T+1)}
\]

\textbf{Combining:}
\[
\left(\frac{1}{2}\right)^{M_i} \leq N \left(\frac{3}{4}\right)^{M_{\mathrm{WM}}}
\]
Taking base-2 logarithms and rearranging:
\[
M_{\mathrm{WM}} \leq \frac{1}{\lg(4/3)}(M_i + \lg N) \approx 2.41\,(M_i + \lg N)
\]

\textbf{Intuition:} Many algorithm mistakes $\Rightarrow$ total weight becomes tiny. A good expert prevents this from happening.

\subsubsection{Tuning the Penalty}

Halving (multiplying by $1/2$ on a mistake) is a special case. More generally, we can multiply by $(1-\varepsilon)$ for any $\varepsilon \in (0,1)$.

\begin{theorem}{Tunable Weighted Majority}{tunable}
With a suitable weighted-majority update using parameter $\varepsilon$:
\[
M_{\mathrm{WM}} \leq 2(1+\varepsilon)\,M_i + O\!\left(\frac{\lg N}{\varepsilon}\right) \quad \text{for every expert } i
\]
\end{theorem}

\textbf{Tradeoff:} Smaller $\varepsilon$ means less overreaction to mistakes (more patient), but slower learning.

% ─────────────────────────────────────────────────────────────────────────────
\subsection{Randomized Experts: The Hedge Algorithm}

\subsubsection{Limitation of Weighted Majority}

Weighted Majority is \textbf{deterministic}. An adversary who can observe our prediction can always choose the opposite outcome, forcing a mistake every round. In fact, there exist sequences where deterministic Weighted Majority makes \emph{about twice} as many mistakes as the best expert. To improve the constant, we stop committing to one answer deterministically.

\subsubsection{Randomized Prediction}

Instead of choosing one expert's advice deterministically, we choose a \textbf{distribution} $\mathbf{p}^{(t)} = (p_1^{(t)}, \ldots, p_N^{(t)})$ over experts. If expert $i$ incurs loss $m_i^{(t)} \in [-1, 1]$ on day $t$, our \textbf{expected loss} is:
\[
\sum_{i=1}^N p_i^{(t)} m_i^{(t)} = \mathbf{p}^{(t)} \cdot \mathbf{m}^{(t)}
\]
This framework covers binary mistakes, but also richer cost structures: money lost, error magnitude, reward minus cost, etc.

\subsubsection{The Hedge Algorithm}

\begin{mdframed}[backgroundcolor=lightblue, linecolor=keyblue, linewidth=1pt]
\textbf{Hedge Algorithm} (parameter $\varepsilon > 0$)

Maintain weights $w_i^{(t)}$. At time $t$:
\begin{enumerate}
  \item Play expert $i$ with probability
    \[
      p_i^{(t)} = \frac{w_i^{(t)}}{\Phi^{(t)}}, \quad \text{where } \Phi^{(t)} = \sum_j w_j^{(t)}
    \]
  \item Observe the loss vector $\mathbf{m}^{(t)}$.
  \item Update weights:
    \[
      w_i^{(t+1)} = w_i^{(t)} \cdot e^{-\varepsilon m_i^{(t)}}
    \]
\end{enumerate}
\end{mdframed}

Low loss $\Rightarrow$ weight increases; high loss $\Rightarrow$ weight decreases. The key distinction from Weighted Majority is the use of the \textbf{exponential update}, which handles continuous losses naturally.

\subsubsection{Hedge Guarantee}

\begin{theorem}{Hedge Regret Bound}{hedge}
If $m_i^{(t)} \in [-1,1]$ and $\varepsilon \leq 1$, then for every expert $i$:
\[
\sum_{t=1}^T \mathbf{p}^{(t)} \cdot \mathbf{m}^{(t)} \;\leq\; \sum_{t=1}^T m_i^{(t)} + \frac{\ln N}{\varepsilon} + \varepsilon T
\]
\end{theorem}

Choosing $\varepsilon \approx \sqrt{\ln N / T}$ balances the two additive terms and gives:
\[
\text{regret} = O\!\left(\sqrt{T \ln N}\right)
\]
The \textbf{average regret} vanishes as $T \to \infty$:
\[
\frac{\text{regret}}{T} = O\!\left(\sqrt{\frac{\ln N}{T}}\right) \;\to\; 0
\]
This is called \textbf{no-regret} or \textbf{Hannan consistency}.

% ─────────────────────────────────────────────────────────────────────────────
\subsection{Why Multiplicative Weights is Everywhere}

The Multiplicative Weights / Hedge framework is a small idea with an enormous footprint across computer science:

\begin{itemize}
  \item \textbf{Boosting} (ML): combine many weak classifiers into a strong one (AdaBoost is a special case).
  \item \textbf{Game theory}: repeated play with Hedge converges to correlated equilibria.
  \item \textbf{Optimization}: fast approximate algorithms for linear programs and semidefinite programs (SDPs).
  \item \textbf{Online learning}: portfolio selection, prediction markets, and adaptive decision making.
  \item \textbf{Algorithm design}: a standard potential-function proof template.
\end{itemize}

\begin{remark}{The recurring pattern}{mw}
\textbf{Multiplicatively reward what worked, but do not completely forget alternatives.}\\[0.3em]
This is the fundamental principle behind multiplicative weights: unlike additive updates that can send weights negative or zero, exponential updates keep all options alive while shifting mass toward better-performing experts.
\end{remark}

% ─────────────────────────────────────────────────────────────────────────────
\subsection{Summary}

\begin{center}
\rowcolors{2}{gray!10}{white}
\begin{tabular}{p{3cm} p{4.5cm} p{5cm}}
\toprule
\textbf{Algorithm} & \textbf{Guarantee} & \textbf{Setting} \\
\midrule
Elimination & $M_{\mathrm{ALG}} \leq \lg N$ & One perfect expert exists \\
Restart Elimination & $M_{\mathrm{ALG}} \leq (M^*+1)\lg N$ & General, but multiplicative overhead \\
Weighted Majority & $M_{\mathrm{WM}} \leq 2.41(M^* + \lg N)$ & General, deterministic \\
Hedge & $\text{regret} = O(\sqrt{T\ln N})$ & General, randomized, continuous losses \\
\bottomrule
\end{tabular}
\end{center}

where $M^* = \min_i M_i$ is the number of mistakes of the best expert.

\textbf{Key takeaways:}
\begin{itemize}
  \item Online algorithms must make irrevocable decisions without seeing the future.
  \item When comparing to full hindsight is too ambitious, \textbf{regret} is the right benchmark.
  \item The \textbf{Elimination algorithm} achieves $\lg N$ mistakes with a perfect expert.
  \item \textbf{Weighted Majority} handles imperfect experts with an additive $O(\lg N)$ overhead, analyzed via a potential function argument.
  \item \textbf{Hedge} uses randomization and exponential updates to achieve sublinear $O(\sqrt{T \ln N})$ regret.
\end{itemize}

\subsection{Introduction: Online vs Offline Algorithms}

\begin{definition}{Offline Algorithm}{offline}
An algorithm where the \textbf{entire input is known in advance}. The algorithm can plan with full information.\\[0.3em]
\textit{Examples:} sorting an array, finding a Minimum Spanning Tree (MST).
\end{definition}

\begin{definition}{Online Algorithm}{online}
An algorithm where \textbf{input arrives one piece at a time}. The algorithm makes \textbf{irrevocable decisions} without knowledge of future requests.\\[0.3em]
\textit{Key principle:} The algorithm must act before it knows the whole story.
\end{definition}

Online decision problems arise in many practical settings:

\begin{center}
\begin{tabular}{ll}
\toprule
\textbf{Setting} & \textbf{Online Decision} \\
\midrule
Cache eviction    & Which page to evict? \\
Online ads        & Which ad to show? \\
Rent or buy       & Keep renting or buy? \\
Expert advice     & Whose advice today? \\
Hiring/selection  & Accept now or wait? \\
\bottomrule
\end{tabular}
\end{center}

A reasonable decision may look poor after the future is revealed. The relevant question is not "was the algorithm psychic?" but rather: \textbf{how badly can it be punished for not being psychic?}

% ─────────────────────────────────────────────────────────────────────────────
\subsection{Measuring Performance}

\begin{definition}{Competitive Ratio}{competitive}
Compare the algorithm's cost with the best solution that knew the whole input in advance (the offline optimum):
\[
\mathrm{ALG}(I) \leq r \cdot \mathrm{OPT}(I) + c
\]
Used when the offline optimum is a meaningful and stable target.
\end{definition}

When comparing to the full offline optimum is too ambitious, we use a weaker but still meaningful benchmark.

\begin{definition}{Regret}{regret}
Compare the algorithm's total loss with the \textbf{best single expert in hindsight}. The extra loss incurred by the algorithm is called \textbf{regret}:
\[
\text{regret} = \mathrm{loss}(\mathrm{ALG}) - \min_i \mathrm{loss}(i)
\]
Used when "best possible" is too strong or too unstable as a benchmark.
\end{definition}

\begin{remark}{Performance benchmarks}{bench}
\textbf{Happiness} $\approx$ outcome $-$ expectation, but happiness is also \textit{relative} to the people next to you. Similarly, algorithm performance is judged relative to what was achievable.
\end{remark}

% ─────────────────────────────────────────────────────────────────────────────
\subsection{Prediction with Expert Advice}

\subsubsection{The Expert Setting}

Consider the following scenario: every morning, $N$ experts each predict YES or NO (e.g., "will it rain?"). We must make our own prediction before observing the outcome.

\textbf{The game.} For $t = 1, \ldots, T$:
\begin{enumerate}
  \item Each expert $i \in \{1, \ldots, N\}$ predicts YES or NO.
  \item The algorithm predicts YES or NO.
  \item The true outcome is revealed.
  \item Everyone who predicted incorrectly receives one \textit{mistake}.
\end{enumerate}

Let $M_{\mathrm{ALG}}$ denote the number of mistakes of the algorithm, and $M_i$ the number of mistakes of expert $i$.

\textbf{Goal:} Keep $M_{\mathrm{ALG}}$ close to $\min_i M_i$.

\subsubsection{Why Compare to the Best Expert?}

\begin{itemize}
  \item We cannot promise to make few mistakes on every sequence (e.g., if all experts say YES and the answer is NO, any algorithm using their advice is in trouble).
  \item But if one expert was consistently good, we would like to discover this quickly.
  \item \textbf{Regret} asks: is $M_{\mathrm{ALG}} - \min_i M_i$ small?
\end{itemize}

% ─────────────────────────────────────────────────────────────────────────────
\subsection{Warmup: One Perfect Expert}

\subsubsection{The Elimination Algorithm}

Suppose at least one expert makes \emph{zero} mistakes. We can exploit this with a simple elimination strategy.

\begin{mdframed}[backgroundcolor=lightblue, linecolor=keyblue, linewidth=1pt]
\textbf{Elimination Algorithm}
\begin{enumerate}
  \item Keep the experts who have made \textbf{no mistakes} so far.
  \item Predict by \textbf{majority vote} among the remaining experts.
  \item After the outcome is revealed, \textbf{delete} every expert who was wrong.
\end{enumerate}
\end{mdframed}

\textbf{Correctness:} The perfect expert is never deleted, so the set of remaining experts is never empty.

\subsubsection{Analysis}

\begin{theorem}{Perfect Expert Guarantee}{perfect}
If at least one expert makes no mistakes, the Elimination algorithm makes \textbf{at most $\lg N$ mistakes}.
\end{theorem}

\textbf{Proof sketch.} Suppose the algorithm makes a mistake. Since it followed the majority of remaining experts, at least half of the remaining experts were wrong — and are now deleted. The pool size decreases:
\[
N \;\to\; \frac{N}{2} \;\to\; \frac{N}{4} \;\to\; \cdots
\]
After more than $\lg N$ mistakes, fewer than one expert would remain — impossible since the perfect expert survives. \hfill$\square$

\textbf{Example:} With $N = 1024$ experts, even a very adversarial sequence causes at most 10 mistakes before the algorithm has essentially identified the perfect expert.

\subsubsection{When Nobody is Perfect}

A first idea: run Elimination until all experts are deleted, then restart with all experts alive. If the best expert makes $M$ mistakes, there are at most $M+1$ restarts, giving:
\[
M_{\mathrm{ALG}} \leq (M+1) \lg N
\]
This bound is often \textbf{too weak}: the overhead grows \emph{multiplicatively} with $M$. We want an \textbf{additive} penalty instead.

% ─────────────────────────────────────────────────────────────────────────────
\subsection{Weighted Majority}

\subsubsection{Soft Elimination}

Instead of deleting an expert after a mistake, we \emph{reduce its influence}:
\begin{itemize}
  \item If correct today: \textbf{keep weight}.
  \item If wrong today: \textbf{halve weight}.
  \item Wrong often: weight \textbf{becomes negligible}.
\end{itemize}
The algorithm listens to a \textbf{weighted majority vote}.

\subsubsection{The Algorithm}

\begin{mdframed}[backgroundcolor=lightblue, linecolor=keyblue, linewidth=1pt]
\textbf{Weighted Majority Algorithm}

Initially set $w_i^{(1)} = 1$ for every expert $i$. For $t = 1, \ldots, T$:
\begin{enumerate}
  \item Each expert predicts YES or NO.
  \item Predict the answer with \textbf{larger total weight}.
  \item After the outcome is revealed, update:
    \[
      w_i^{(t+1)} = \begin{cases} w_i^{(t)} / 2 & \text{if expert } i \text{ was wrong,} \\ w_i^{(t)} & \text{otherwise.} \end{cases}
    \]
\end{enumerate}
\end{mdframed}

Weights act as \textbf{trust scores}: they only decrease when the expert makes a mistake, so a good expert retains much more weight than a bad one.

\subsubsection{Guarantee}

\begin{theorem}{Weighted Majority Bound}{wm}
Let $M_{\mathrm{WM}}$ be the number of mistakes of Weighted Majority. For every expert $i$:
\[
M_{\mathrm{WM}} \leq 2.41 \cdot (M_i + \lg N)
\]
Therefore:
\[
M_{\mathrm{WM}} \leq 2.41 \cdot \left(\min_i M_i + \lg N\right)
\]
\end{theorem}

This is much better than the restarting bound when $M_i$ is large: the overhead is now \textbf{additive} in $\lg N$.

\subsubsection{Analysis via Potential Function}

Define the \textbf{potential} (total weight):
\[
\Phi^{(t)} = \sum_{i=1}^{N} w_i^{(t)}
\]

\textbf{Upper bound} (from algorithm's mistakes): When the algorithm makes a mistake, it followed the weighted majority, so at least half the total weight was on wrong experts. Those weights are halved, giving:
\[
\Phi^{(t+1)} \leq \frac{1}{2}\Phi^{(t)} + \frac{1}{2} \cdot \frac{1}{2}\Phi^{(t)} = \frac{3}{4}\Phi^{(t)}
\]
After $M_{\mathrm{WM}}$ mistakes:
\[
\Phi^{(T+1)} \leq N \left(\frac{3}{4}\right)^{M_{\mathrm{WM}}}
\]

\textbf{Lower bound} (from a good expert): If expert $i$ makes $M_i$ mistakes, its final weight is $(1/2)^{M_i}$. Since the total weight is at least this:
\[
\left(\frac{1}{2}\right)^{M_i} \leq \Phi^{(T+1)}
\]

\textbf{Combining:}
\[
\left(\frac{1}{2}\right)^{M_i} \leq N \left(\frac{3}{4}\right)^{M_{\mathrm{WM}}}
\]
Taking base-2 logarithms and rearranging:
\[
M_{\mathrm{WM}} \leq \frac{1}{\lg(4/3)}(M_i + \lg N) \approx 2.41\,(M_i + \lg N)
\]

\textbf{Intuition:} Many algorithm mistakes $\Rightarrow$ total weight becomes tiny. A good expert prevents this from happening.

\subsubsection{Tuning the Penalty}

Halving (multiplying by $1/2$ on a mistake) is a special case. More generally, we can multiply by $(1-\varepsilon)$ for any $\varepsilon \in (0,1)$.

\begin{theorem}{Tunable Weighted Majority}{tunable}
With a suitable weighted-majority update using parameter $\varepsilon$:
\[
M_{\mathrm{WM}} \leq 2(1+\varepsilon)\,M_i + O\!\left(\frac{\lg N}{\varepsilon}\right) \quad \text{for every expert } i
\]
\end{theorem}

\textbf{Tradeoff:} Smaller $\varepsilon$ means less overreaction to mistakes (more patient), but slower learning.

% ─────────────────────────────────────────────────────────────────────────────
\subsection{Randomized Experts: The Hedge Algorithm}

\subsubsection{Limitation of Weighted Majority}

Weighted Majority is \textbf{deterministic}. An adversary who can observe our prediction can always choose the opposite outcome, forcing a mistake every round. In fact, there exist sequences where deterministic Weighted Majority makes \emph{about twice} as many mistakes as the best expert. To improve the constant, we stop committing to one answer deterministically.

\subsubsection{Randomized Prediction}

Instead of choosing one expert's advice deterministically, we choose a \textbf{distribution} $\mathbf{p}^{(t)} = (p_1^{(t)}, \ldots, p_N^{(t)})$ over experts. If expert $i$ incurs loss $m_i^{(t)} \in [-1, 1]$ on day $t$, our \textbf{expected loss} is:
\[
\sum_{i=1}^N p_i^{(t)} m_i^{(t)} = \mathbf{p}^{(t)} \cdot \mathbf{m}^{(t)}
\]
This framework covers binary mistakes, but also richer cost structures: money lost, error magnitude, reward minus cost, etc.

\subsubsection{The Hedge Algorithm}

\begin{mdframed}[backgroundcolor=lightblue, linecolor=keyblue, linewidth=1pt]
\textbf{Hedge Algorithm} (parameter $\varepsilon > 0$)

Maintain weights $w_i^{(t)}$. At time $t$:
\begin{enumerate}
  \item Play expert $i$ with probability
    \[
      p_i^{(t)} = \frac{w_i^{(t)}}{\Phi^{(t)}}, \quad \text{where } \Phi^{(t)} = \sum_j w_j^{(t)}
    \]
  \item Observe the loss vector $\mathbf{m}^{(t)}$.
  \item Update weights:
    \[
      w_i^{(t+1)} = w_i^{(t)} \cdot e^{-\varepsilon m_i^{(t)}}
    \]
\end{enumerate}
\end{mdframed}

Low loss $\Rightarrow$ weight increases; high loss $\Rightarrow$ weight decreases. The key distinction from Weighted Majority is the use of the \textbf{exponential update}, which handles continuous losses naturally.

\subsubsection{Hedge Guarantee}

\begin{theorem}{Hedge Regret Bound}{hedge}
If $m_i^{(t)} \in [-1,1]$ and $\varepsilon \leq 1$, then for every expert $i$:
\[
\sum_{t=1}^T \mathbf{p}^{(t)} \cdot \mathbf{m}^{(t)} \;\leq\; \sum_{t=1}^T m_i^{(t)} + \frac{\ln N}{\varepsilon} + \varepsilon T
\]
\end{theorem}

Choosing $\varepsilon \approx \sqrt{\ln N / T}$ balances the two additive terms and gives:
\[
\text{regret} = O\!\left(\sqrt{T \ln N}\right)
\]
The \textbf{average regret} vanishes as $T \to \infty$:
\[
\frac{\text{regret}}{T} = O\!\left(\sqrt{\frac{\ln N}{T}}\right) \;\to\; 0
\]
This is called \textbf{no-regret} or \textbf{Hannan consistency}.

% ─────────────────────────────────────────────────────────────────────────────
\subsection{Why Multiplicative Weights is Everywhere}

The Multiplicative Weights / Hedge framework is a small idea with an enormous footprint across computer science:

\begin{itemize}
  \item \textbf{Boosting} (ML): combine many weak classifiers into a strong one (AdaBoost is a special case).
  \item \textbf{Game theory}: repeated play with Hedge converges to correlated equilibria.
  \item \textbf{Optimization}: fast approximate algorithms for linear programs and semidefinite programs (SDPs).
  \item \textbf{Online learning}: portfolio selection, prediction markets, and adaptive decision making.
  \item \textbf{Algorithm design}: a standard potential-function proof template.
\end{itemize}

\begin{remark}{The recurring pattern}{mw}
\textbf{Multiplicatively reward what worked, but do not completely forget alternatives.}\\[0.3em]
This is the fundamental principle behind multiplicative weights: unlike additive updates that can send weights negative or zero, exponential updates keep all options alive while shifting mass toward better-performing experts.
\end{remark}

% ─────────────────────────────────────────────────────────────────────────────
\subsection{Summary}

\begin{center}
\rowcolors{2}{gray!10}{white}
\begin{tabular}{p{3cm} p{4.5cm} p{5cm}}
\toprule
\textbf{Algorithm} & \textbf{Guarantee} & \textbf{Setting} \\
\midrule
Elimination & $M_{\mathrm{ALG}} \leq \lg N$ & One perfect expert exists \\
Restart Elimination & $M_{\mathrm{ALG}} \leq (M^*+1)\lg N$ & General, but multiplicative overhead \\
Weighted Majority & $M_{\mathrm{WM}} \leq 2.41(M^* + \lg N)$ & General, deterministic \\
Hedge & $\text{regret} = O(\sqrt{T\ln N})$ & General, randomized, continuous losses \\
\bottomrule
\end{tabular}
\end{center}

where $M^* = \min_i M_i$ is the number of mistakes of the best expert.

\textbf{Key takeaways:}
\begin{itemize}
  \item Online algorithms must make irrevocable decisions without seeing the future.
  \item When comparing to full hindsight is too ambitious, \textbf{regret} is the right benchmark.
  \item The \textbf{Elimination algorithm} achieves $\lg N$ mistakes with a perfect expert.
  \item \textbf{Weighted Majority} handles imperfect experts with an additive $O(\lg N)$ overhead, analyzed via a potential function argument.
  \item \textbf{Hedge} uses randomization and exponential updates to achieve sublinear $O(\sqrt{T \ln N})$ regret.
\end{itemize}